\section{Interactive computational laboratory}\label{sec:lab}

\begin{figure*}
\includegraphics[width=\textwidth]{figures/fig_dashboard.pdf}
\caption{Snapshot of the interactive laboratory for the chaotic reference state ($\mu=-0.8$, $\gamma=0.25$, $A=0.8$, $\Omega=2/3$). (a) Schematic with the fading motion trail of the mass, drawn through the Hermite interpolant. (b) Energy landscape with the separatrix energy and the current energy. (c) Phase portrait: logarithmic occupation histogram, separatrix of the unforced system (black), recent trajectory, and equilibria of the unforced system. (d) Stroboscopic Poincar\'e section. (e) Normalized spring force $F_s/(kL)$ (blue) and collar height $y_c/L$ (orange). (f) Power spectral density of the spring force as a function of frequency in units of the drive frequency $f_d=\Omega/(2\pi)$.}
\label{fig:dashboard}
\end{figure*}

Every result of this article is produced by an open laboratory written in Python. The package \texttt{piar} contains one module per layer of the analysis: \texttt{physics} (the model, its closed-form invariants, and compiled right-hand sides), \texttt{integrators} (a compiled DOP853, velocity Verlet, fourth-order Runge--Kutta, and exact periods), \texttt{chaos} (stroboscopic maps, continuation, Lyapunov spectra, Floquet multipliers, Feigenbaum cascades, Melnikov integrals, and correlation sums), \texttt{stochastic} (BAOAB Langevin dynamics), \texttt{embedding} (average mutual information and false nearest neighbors), \texttt{fprm} (conditioning vectors, conditional flows, cVAE, DDPM, Gaussian processes, physics residuals, and metrics), \texttt{smooth} (rendering of trajectories), and \texttt{dashboard}. Array operations rely on NumPy~\cite{Harris2020} and SciPy~\cite{Virtanen2020}, numerical kernels are compiled with Numba~\cite{Lam2015}, the neural networks are implemented in PyTorch~\cite{Paszke2019}, the flows are provided by \texttt{zuko}~\cite{zuko}, the Gaussian processes by scikit-learn~\cite{Pedregosa2011}, and the figures are drawn with Matplotlib~\cite{Hunter2007}.

Trajectories are rendered from continuous representations of the solution. Conservative trajectories are drawn from the dense output of DOP853 (Sec.~\ref{sec:cases_method}). For fixed-step integration, the integrator provides $\theta$ and $\dot\theta$ at each grid point and the equation of motion provides $\ddot\theta$ at the cost of one evaluation per stored point, so that each interval is filled with the cubic Hermite interpolants of $\theta$ and $\dot\theta$, which are continuously differentiable and accurate to fourth order in the step~\cite{HairerNW1993}; this is how the driven trajectories of Figs.~\ref{fig:attractor}(a) and \ref{fig:reconstruction} are drawn. Angles are wrapped onto the cylinder with explicit breaks at the cut $\theta=\pm\pi$, so that rotations are not crossed by spurious horizontal segments.

The interactive part of the laboratory follows the environments long used in physics teaching, such as Easy Java Simulations~\cite{Esquembre2004} and the Open Source Physics project~\cite{Christian2011}. The interactive dashboard (Fig.~\ref{fig:dashboard}) is a Jupyter notebook in which sliders for $(\mu,\gamma,A,\Omega)$ and for the initial condition drive six linked panels: a schematic with a fading motion trail, the energy landscape with the current energy, the phase portrait, the stroboscopic Poincar\'e section, the two measurable observables, and the power spectrum of the spring force. The notebook also draws the four diagnostics of Fig.~\ref{fig:cases} for any case and animates the pendulum and its phase-space trajectory from the same Hermite interpolant. A browser-based version, written in JavaScript and requiring no installation, reproduces the five cases of Table~\ref{tab:cases} and the chaotic reference state; it integrates the equation of motion with fourth-order Runge--Kutta at a step of $2\times10^{-3}\,t_c$, keeps the energy of the conservative cases constant to better than $10^{-10}\,MgL$ over 100~s, stores every step in the motion trail, and updates the four diagnostics in real time.

All figures follow a common style: Computer Modern typography rendered by \LaTeX, axes without top and right spines, and perceptually uniform scientific colormaps that remain legible under color-vision deficiency (\textit{batlow} for magnitudes and \textit{vik} for signed quantities~\cite{Crameri2020}), together with a categorical palette whose adjacent pairs were checked for separability under simulated color-vision deficiency. Eight experiment scripts and five figure scripts regenerate every figure and number from fixed random seeds and write the numerical results to JSON files. Unit tests check the closed-form invariants, the agreement of the compiled DOP853 with SciPy, the second-order convergence of velocity Verlet, the exact periods, the determinant of the Poincar\'e map, the Lyapunov sum rule, the pendulum limit of the Melnikov integrals, and the accuracy of the Hermite rendering.

\paragraph*{Use of generative artificial intelligence.} A generative artificial-intelligence tool, Claude (model Claude Opus 5.5; Anthropic PBC, San Francisco, CA, USA), accessed through the Claude application in September 2026, was used in three parts of this work: (i) writing and debugging the analysis code, namely the \texttt{piar} package, the experiment and figure scripts, the unit tests, the notebook, and the browser-based simulation; (ii) running the numerical experiments and documenting their settings and outputs; and (iii) drafting and editing the manuscript. The tool was used to accelerate software development, testing, and documentation. Because this use could affect the findings, every figure and number reported in this article is regenerated by the released code from fixed random seeds, the unit tests check the closed-form invariants, the integrators, and the chaos diagnostics against analytical or independent results, and the bibliographic data of every reference were checked against publisher or index records. The author takes full responsibility for the code, the numerical results, and the text.
